\documentclass[12pt,a4paper]{article}

\usepackage[utf8]{inputenc}
\usepackage[T1]{fontenc}
\usepackage{geometry}
\usepackage{hyperref}
\usepackage{setspace}
\usepackage{parskip}
\usepackage{booktabs}
\usepackage{array}
\usepackage{longtable}
\usepackage{xcolor}
\usepackage{soul}
\usepackage{amsmath}
\usepackage{amssymb}

\geometry{margin=2.5cm}
\onehalfspacing

% Helper for Giray placeholders
\newcommand{\giray}[1]{\textbf{[GIRAY: #1]}}

\title{\textbf{CAST+: Augmenting the ACLED Conflict Alert System }}
\author{Rosa Daneshmandnia \and Giray \"{U}nl\"{u}}
\date{May 2026}

\begin{document}

\maketitle

% ============================================================
\begin{abstract}
The ACLED Conflict Alert System (CAST) is a global prediction tool that
forecasts monthly political violence at the subnational level. While
CAST provides strong baseline forecasts, it tends to extend recent
patterns forward, limiting its ability to detect meaningful shifts in
conflict dynamics. This report presents CAST+, an augmented system
developed in collaboration with ACLED and Algorithmic Governance over
six weeks in early 2026. CAST+ makes two complementary contributions.
First, it expands the input feature space with new open-source data
signals including news-scraped risk indicators, World Bank macroeconomic
variables, public holiday calendars, and religious composition data,
alongside engineered interaction features. Second, it introduces a
two-stage CatBoost architecture that directly addresses the zero-inflated
structure of monthly conflict change, predicting whether change will
occur before estimating its magnitude. Evaluated on a six-month
chronological holdout across 2,205 global Admin-1 regions, the best
CAST+ configuration achieves a mean absolute error of 1.589, representing
a 24.6\% improvement over the CAST baseline of 2.344, surpassing the
project target of 20\%. On the top-10 most active regions, CAST+ achieves
a 23.1\% improvement over the baseline, outperforming all alternative
configurations tested. A directional accuracy of 79 to 82\% across
targets confirms the model is learning meaningful escalation patterns
rather than simply copying recent conflict levels. All experiments are
fully reproducible via the project repository at
\url{https://github.com/rapsoj/predicting-global-conflict}.
\end{abstract}

\tableofcontents
\newpage

% ============================================================
\section{Introduction}

Political violence, encompassing armed conflict, targeted attacks on
civilians, and the use of explosive and remote weaponry, remains one
of the most consequential and difficult-to-anticipate threats to human
security globally. The ability to forecast where and when violence is
likely to escalate is of significant value to policymakers, humanitarian
organisations, and conflict researchers, who depend on early warning
signals to allocate resources, plan responses, and where possible,
prevent harm before it occurs.

The Armed Conflict Location and Event Data Project (ACLED) is the
leading global repository of disaggregated conflict data, recording
political violence and protest events across every country in the world
with no minimum fatality threshold for inclusion. Built on this
foundation, ACLED's Conflict Alert System (CAST) is a global prediction
tool that forecasts the number of political violence events expected in
the coming six months, at both the country level and the first
administrative division (Admin-1) level. CAST disaggregates its forecasts
across three event types, namely battles, explosions and remote violence,
and violence against civilians, providing users with granular, actionable
predictions rather than a single aggregate conflict index.

Despite its scope and utility, CAST exhibits a characteristic limitation
common to many conflict forecasting systems: it behaves largely as a
persistence model, extrapolating recent patterns forward in time. This
design produces reliable forecasts in stable conflict environments but
systematically underperforms in situations of meaningful change, precisely
the escalation dynamics, emerging hotspots, and shifting conflict
trajectories that early warning systems are most needed to detect.

This project, conducted over approximately six weeks in early 2026, aims
to address this limitation through two complementary contributions. First,
the feature space is expanded by integrating new open-source data signals
including news-scraped risk indicators, World Bank macroeconomic
indicators, public holiday calendars, religious composition data, and a
set of engineered interaction features. Second, a systematic model
benchmarking study is conducted across multiple architectures and an
extended set of engineered features, identifying the configurations that
most improve forecast accuracy and change detection relative to the
repository baseline.

The result is CAST+: an augmented forecasting pipeline that retains the
interpretability and global coverage of the original system while
improving its sensitivity to conflict dynamics that deviate from
established patterns.


% ============================================================
\section{Background and Related Work}

\subsection{The ACLED Conflict Alert System}

% FLAG: After Kai meeting, confirm whether Section 2.1 should describe
% the repository baseline (RF) or the current live CAST (LGBM-Tweedie).
% Currently framed around the repository baseline used in this project.

CAST structures its predictions at the Admin-1 month level, training
separate models for each country on data from 2018 onward. For each
country and each of the three target event types, the repository baseline
fits and evaluates four candidate model configurations, namely a random
forest and an XGBoost model, each run against a simple and a complex
predictor list, using four-fold time-series cross-validation with a
six-month holdout window. The best-performing configuration for each
country is selected and retrained on the full available data to generate
forecasts for the following six months.

The predictor set used by CAST is derived primarily from ACLED data
itself: lagged event counts from the previous month, spillover from
neighbouring administrative divisions, actor interaction types, fatality
rates, and strategic development events such as peace agreements. This is
supplemented by external sources including population estimates from
WorldPop and infant mortality rates from the Center for International
Earth Science Information Network (CIESIN) as a proxy for economic
development. Forecasts are produced at the Admin-1 level and aggregated
upward to the country and global level.

\subsection{Limitations of the Existing Approach}

The persistence-like behaviour of the baseline follows naturally from its
feature construction: when the strongest predictors are lagged event
counts from the previous month, the model tends to forecast next month's
violence as a function of last month's. This is a well-documented
challenge in conflict forecasting. Models trained on historical conflict
data tend to be well-calibrated for ongoing, stable conflicts but
struggle with onset and escalation events, which are low-frequency,
high-impact shifts that are underrepresented in the training distribution.

The CAST methodology acknowledges this asymmetry directly, noting that
forecast accuracy is not consistent across all countries or administrative
divisions, and that a stable conflict with an established pattern of
events produces more accurate forecasts than a new, highly volatile
conflict with no established pattern. Improving performance on the latter
category is the central motivation for CAST+.

Additionally, 86 to 92\% of monthly delta values across all Admin-1
regions are zero, meaning conflict levels do not change from one month
to the next. A single-stage regression model wastes capacity predicting
these zero-change months rather than concentrating on the months where
change actually occurs.

\subsection{Project Objectives}

This project sets out to improve the model's ability to detect meaningful
changes in conflict risk over a six-month horizon, with a target of 20\%
improvement over the baseline. Concretely, this involves three objectives:
(1) expanding the input feature space with new data sources that carry
signal about changing conflict conditions; (2) engineering features that
better capture temporal dynamics and structural vulnerability; and (3)
identifying model architectures that are better suited to the
zero-inflated, count-based nature of the prediction task than the
random forest baseline.


% ============================================================
\section{Data and Datasets}

The CAST+ pipeline integrates five data sources. The primary source is
the ACLED event-level dataset. Four additional open-source datasets were
integrated to enrich the feature space.

\subsection{ACLED Event Data}

The primary data source is the ACLED raw event file, covering
1 January 1997 to 20 June 2025, comprising 2,619,146 raw events. The
pipeline filters to 2018 onward (2,241,762 events, 85.6\% of the raw
file) before monthly aggregation, consistent with the CAST methodology.
Pre-2018 data is excluded due to inconsistent global coverage in ACLED's
historical collection. After monthly aggregation, the dataset spans
approximately 278,821 region-month rows across approximately 3,097 global
Admin-1 regions. Admin-1 region boundaries are derived from the Natural
Earth Admin-1 shapefiles.

\subsection{News-Scraped Risk Indicators}

News coverage of seven conflict-relevant risk categories was collected
via a four-stage pipeline implemented in \texttt{src/scraping/}.
The \texttt{GNewsFetcher} module queries the Google News RSS feed for
each country--metric--year combination using per-metric keyword
templates (e.g.\ ``crop failure OR harvest crisis'' for crop failure).
Full article text is retrieved with a headless Chromium browser
(Playwright, up to 75 concurrent tabs). A fuzzy keyword pre-filter
removes off-topic articles; surviving articles are sent to an OpenAI
language model that assigns each to one of the seven categories and
records the country, metric, and month-year. Results are aggregated to
country-month counts in \texttt{data/raw/master\_raw.csv}: 26,530
records covering 236 countries from January 2018 to December 2024.
The seven categories are: contested election, crop failure, economic
concern, ethnic tension, military coup, natural disaster, and political
assassination. Coverage is sparse by design---between 5 and 21\% of
region-months have a non-zero value for any given indicator.

\subsection{World Bank Macroeconomic Indicators}

Three macroeconomic indicators were retrieved via the World Bank API
and integrated as country-level annual features:

\begin{itemize}
    \item \textbf{Inflation} (CPI annual percentage change,
          indicator FP.CPI.TOTL.ZG): captures economic shock conditions.
          Analysis at the country-year level shows that inflation does not
          exhibit a strong linear relationship with conflict counts, but
          high inflation years (top quintile) are associated with
          substantially higher violence, particularly in low income
          countries. Inflation functions as a dynamic accelerator of
          conflict in economically vulnerable contexts.
    \item \textbf{Youth unemployment} (SL.UEM.NEET.ZS): captures
          mobilisation capacity. Moderate positive correlation with
          conflict totals (0.20 to 0.33) suggests youth unemployment
          contributes to the pool of potential recruits into armed groups.
    \item \textbf{Income inequality} (Gini coefficient,
          SI.POV.GINI): the strongest structural signal in this dataset.
          Country-year analysis reveals a threshold-like relationship:
          high inequality countries exhibit dramatically higher violence
          against civilians (mean 535 events per year) compared to low
          inequality countries (mean 12 events per year). This separation
          persists after log transformation and within World Bank income
          tiers, indicating inequality is not merely a proxy for
          development level.
\end{itemize}

In addition, \textbf{income level classification} (Low, Lower-middle,
Upper-middle, High income) was retrieved from the World Bank country
metadata API and encoded as an ordinal variable (LIC=1 through HIC=4),
providing a structural institutional buffer signal.

\textbf{Leakage prevention:} World Bank data is published with a 12 to
18 month lag. All three continuous indicators are shifted by one prior
year before joining to the model dataset. Income level is a structural
classification and requires no lag. Sparse indicators such as the Gini
coefficient are filled using within-country backward and forward fill
before broadcasting annual values to all 12 months of each year.

\subsection{Public Holiday Calendar}

A public holiday dataset covering 233 countries was integrated, comprising
89,217 national holiday records after cleaning. Holiday names were
standardised across 34 spelling variants and tagged by religious
tradition, producing eight monthly count features per country:
\texttt{christian\_holiday\_count}, 
\texttt{islam\_holiday\_count},
\texttt{shia\_holiday\_count}, \texttt{hindu\_holiday\_count},
\texttt{buddhist\_holiday\_count}, \texttt{jewish\_holiday\_count},
\texttt{cultural\_holiday\_count}, and
\texttt{nonreligious\_holiday\_count}.

Unlike other features in the pipeline, holidays are deterministic
calendar facts known in advance. They are therefore used at the
current month without a lag, which introduces no data leakage while
preserving the full information content.

\subsection{Religious Composition Data}

Religious composition data was sourced from the World Religion Project
(WRP) national dataset, joined to ISO3 country codes via the Correlates
of War (COW) country code mapping. The 2010 snapshot was used as a
static structural baseline capturing the top three religious denominations
by population share for each country:
\texttt{majority\_religion}, \texttt{majority\_pct},
\texttt{minority1\_religion}, \texttt{minority1\_pct},
\texttt{minority2\_religion}, \texttt{minority2\_pct}, and
\texttt{nonreligpct}.

The motivation for including religious composition is that holidays only
become meaningful conflict signals when considered in their religious
context. An Islamic holiday in a country with a Muslim majority carries
different social significance than the same date in a country where
Muslims are a small minority. This dataset enables the holiday features
to be contextualised by the religious structure of each country.

\subsection{Data Merging and Integration}

All five datasets are joined to a single Admin-1 monthly panel indexed
by (\texttt{matched\_admin1\_id}, \texttt{month\_year}). Country-level
features (World Bank, religion, risk indicators, holidays) are broadcast
to all Admin-1 regions within the same country via an ISO3 country code
extracted from the \texttt{matched\_admin1\_id} prefix. Risk indicator counts are matched to the panel on (ISO3 country code,
month-year) and broadcast to all Admin-1 regions within that country,
then lagged one month ($t-1$) to prevent data leakage.

The full merged and engineered dataset spans 289,890 rows across
3,221 Admin-1 regions and 78 feature columns.


% ============================================================
\section{Feature Engineering}

Features are organised into cumulative layers. The first two layers
constitute the CAST baseline predictor set. Subsequent layers represent
enrichments developed in this project.

\subsection{Layer 1: Autoregressive Lag Features (ACLED)}

One-month lagged counts for all 8 ACLED event types: Battles, Explosions
and Remote violence, Protests, Riots, Strategic developments, Violence
against civilians, Excessive force against protesters, and Agreement.
Additionally, lag-1 of event counts summed over all bordering Admin-1
regions is included as a spatial spillover signal. Neighbour adjacency
is derived from the Natural Earth Admin-1 shapefile by polygon
intersection. All lags are computed per region to prevent cross-region
leakage.

\subsection{Layer 2: Temporal Features}

A set of time trend variables capturing seasonal and trend components:
a linear month trend since January 2018, year, month dummies, and quarter
dummies. These features capture unobserved time-varying patterns including
seasonality and long-run conflict trends.

\subsection{Layer 3: News-Scraped Risk Indicators}

Seven country-month risk indicators derived from the pipeline described
in Section~3.2 are included as features:
\texttt{risk\_contested\_election},
\texttt{risk\_crop\_failure},
\texttt{risk\_economic\_concern},
\texttt{risk\_ethnic\_tension},
\texttt{risk\_military\_coup},
\texttt{risk\_natural\_disaster}, and
\texttt{risk\_political\_assassination}.
Each records the monthly article count for that risk category,
matched via ISO3 country code and broadcast to all Admin-1 regions
within that country. A one-month lag ($t-1$) is applied to prevent
data leakage. As shown in Table~\ref{tab:ablation_features}, these
features add noise in practice and are candidates for removal in a
production pipeline.

\subsection{Layer 4: World Bank Macroeconomic Features}

The four World Bank features (inflation, youth unemployment, income
inequality, income level code) described in Section 3.3 are added as
direct inputs. Prior-year shift is applied to the three continuous
indicators to prevent data leakage, as described above.

\subsection{Layer 5: Engineered World Bank Features}

Beyond the raw World Bank indicators, ten derived features were
constructed to capture non-linear interactions and threshold effects
identified in the data analysis:

\begin{itemize}
    \item \textbf{Structural vulnerability}: income inequality multiplied
          by youth unemployment, representing the combination of
          chronic grievance and a mobilisable population.
    \item \textbf{Inflation shock}: binary flag for top-decile inflation
          years, capturing the non-linear shock threshold identified
          in the country-year analysis.
    \item \textbf{Shock vulnerability}: inflation shock multiplied by
          structural vulnerability, combining the acute shock with the
          underlying structural condition.
    \item \textbf{Capacity adjusted risk}: structural vulnerability
          divided by income level, representing the institutional
          buffering effect of higher income on instability.
    \item \textbf{Inflation change}: first difference of lagged
          inflation, capturing acceleration rather than level.
    \item \textbf{High inequality flag}: binary flag for top-quartile
          inequality, capturing the threshold effect observed in the
          country-year analysis.
    \item \textbf{Inflation shock in poor countries}: interaction of
          inflation shock and low income flag, isolating the most
          vulnerable context.
    \item \textbf{Conflict momentum}: sum of Battles and Violence
          against civilians plus their neighbour aggregates, capturing
          overall conflict pressure.
    \item \textbf{Violence escalation}: ratio of violent events
          (Battles plus Violence against civilians) to non-violent
          activity (Protests plus Riots plus 1), capturing the
          shift from political mobilisation to armed violence.
    \item \textbf{Macro conflict pressure}: structural vulnerability
          multiplied by conflict momentum, combining the structural
          enabling condition with current conflict intensity.
\end{itemize}

\subsection{Layer 6: Holiday Features}

The eight religion-specific holiday count columns described in Section
3.4 are added as current-month features. No lag is applied, as holidays
are known calendar facts.

\subsection{Layer 7: Holiday and Religion Interaction Features}

To make holiday features meaningful in context, three layers of
holiday-religion interaction features were constructed:

\begin{itemize}
    \item \textbf{Relevance-weighted holiday counts}: each holiday type
          is weighted by the corresponding religion's population share
          in that country. For example, Islam holiday count is multiplied
          by the majority population share when the majority religion is
          Islam. This produces \texttt{weighted\_holidays\_majority},
          \texttt{weighted\_holidays\_minority1}, and
          \texttt{weighted\_holidays\_minority2}.
    \item \textbf{Minority tension signal}: a minority religion has a
          holiday but is not the dominant group. This is a proxy for
          potential inter-group mobilisation or friction:
          \texttt{minority\_tension\_minority1},
          \texttt{minority\_tension\_minority2}, and
          \texttt{minority\_tension\_total}.
    \item \textbf{Total religious mobilisation}: composite sum of all
          weighted holiday signals across majority and minority groups.
\end{itemize}

All seven interaction features are lagged by one month (t-1) to maintain
consistency with the pipeline lag policy for dynamic features.

\subsection{Layer 8: Country-Level Hierarchy Features}

Four leave-one-out (LOO) country aggregate features provide a
national-context signal without leaking the focal region's own
activity: \texttt{country\_battles\_excl~$(t-1)$},
\texttt{country\_remote\_excl~$(t-1)$},
\texttt{country\_vac\_excl~$(t-1)$}, and
\texttt{country\_total\_excl~$(t-1)$}. Each is the sum of the
corresponding conflict type across all Admin-1 regions in the same
country, excluding the focal region. The LOO exclusion prevents
multicollinearity with the Layer~1 autoregressive features; the
$t-1$ lag prevents leakage since the current-month country total
contains the focal region's own current-month count.

In the LGBM-Tweedie V2 ablation, the +Country tier delivers the
single largest improvement: $-4.54$ MAE ($-7.9\%$), more than three
times the combined contribution of all other enrichment layers.
This benefit does not transfer to the two-stage CatBoost framework
($<1\%$ gain), where the Stage~1 onset classifier already captures
national-context information from the data directly.

\subsection{Layer 9: Engineered Cross-Variable Features}

Nine cross-variable features were constructed from the Layer~1
lag-1 counts: lag-2 values for each of the three targets (enabling
implicit slope computation); an \texttt{organized\_violence~$(t-1)$}
aggregate (Battles + Remote + VaC); a binary \texttt{is\_active~$(t-1)$}
zero-inflation flag; a \texttt{battles\_x\_remote~$(t-1)$}
co-escalation interaction term; and three-month rolling averages for
each target. A further nine trajectory features (first-difference
slope, second-difference acceleration, and neighbour slope per target)
provide a marginal additional improvement of $-0.25$ MAE when added
to the V2 stack.

\subsection{Feature Ablation Results}

A feature ablation study was conducted on 100 regions to measure the
marginal contribution of each feature group to the two-stage CatBoost
model. The results are summarised in Table~\ref{tab:ablation_features}.

\begin{table}[h]
\centering
\caption{Feature ablation results on 100 regions. Negative percentage
indicates the removed group improves MAE (adds noise). Positive
indicates the group contributes positively.}
\label{tab:ablation_features}
\begin{tabular}{lrrrl}
\toprule
Removed group & Battles & Remote Violence & VAC & Conclusion \\
\midrule
No Religion    & 0.0\%   & 0.0\%   & 0.0\%  & No contribution \\
No Holidays    & -2.4\%  & 0.0\%   & +1.4\% & Mixed, marginal \\
No World Bank  & -5.1\%  & +0.7\%  & +0.5\% & Slight noise for Battles \\
No Risk        & -8.4\%  & -2.8\%  & -3.8\% & Adds noise \\
\bottomrule
\end{tabular}
\end{table}

Key findings from the ablation are as follows. Religion features make no
measurable contribution to the model. Holiday features have a marginal
mixed effect, helping VAC slightly while adding minor noise for Battles.
World Bank features add slight noise for Battles but provide marginal
signal for Remote violence and VAC. Risk features add noise across all
three targets, with removing them improving MAE for every target. Despite
being theoretically meaningful (crop failure, ethnic tension, contested
elections), risk features do not help the model in practice and are
candidates for removal in a production pipeline.

The model is primarily driven by recent conflict history. Stage 2
(magnitude prediction) is dominated by \texttt{violence\_escalation}
(importance 19.5) and conflict lags. Stage 1 (change classifier) shows
a more distributed importance profile, with \texttt{linear\_month\_trend},
\texttt{conflict\_momentum}, and \texttt{violence\_escalation} leading,
but risk features, economic indicators, and religious mobilisation all
contributing meaningfully to detecting whether change will occur.


% ============================================================
\section{Methodology}

\subsection{Problem Formulation}

The forecasting task is defined at the Admin-1 region level as a monthly
prediction of three target event counts: Battles, Explosions and Remote
violence, and Violence against civilians. Each target is treated as an
independent regression problem.

A key design decision in CAST+ is the \textbf{delta formulation}: instead
of predicting raw event counts $Y(t)$, the model predicts the
month-to-month change $\Delta Y = Y(t) - Y(t-1)$. Final count predictions
are recovered at inference via:
\[
    \hat{Y}(t) = Y(t-1) + \hat{\Delta Y}
\]

This formulation is motivated by the structure of the data: 86 to 92\%
of monthly delta values are zero, meaning conflict levels do not change
from one month to the next. Predicting change is structurally simpler
because the model only needs to explain deviations from the current
state rather than reconstruct absolute levels. It also directly aligns
the model's objective with the project's stated goal of change detection.

\subsection{Evaluation Design}

\textbf{Primary evaluation:} All 2,205 valid Admin-1 regions globally,
using a six-month chronological holdout. Regions with fewer than 12
months of data or no variation in delta targets are excluded. This
evaluation is the primary reported result because CAST is a global
forecasting tool and should be evaluated on a globally representative
population.

\textbf{Secondary evaluation:} The top-10 most active regions by total
event count during the holdout period, for direct comparison with the
alternative ablation experiments. These regions are: UKR-Donetsk,
UKR-Sumy, RUS-Belgorod, UKR-Kherson, PSX-Gaza, UKR-Kharkiv,
UKR-Zaporizhzhia, RUS-Kursk, UKR-Chernihiv, and PSX-West Bank.

\textbf{Train-test split:} Strictly chronological. Training data uses
all months up to $T-6$. Test data uses the last 6 months. No shuffling
is applied.

\textbf{Primary metric:} Mean Absolute Error (MAE) at the count level.
Delta predictions are converted back to counts before computing MAE,
ensuring fair comparison with the CAST baseline which predicts raw counts.

\textbf{Secondary metric:} Directional accuracy, measuring whether the
model correctly predicts the direction of change (escalating, stable, or
de-escalating). This metric is more meaningful for an early warning tool
than raw MAE, as it evaluates whether the model correctly signals the
onset of escalation rather than its exact magnitude.

\textbf{Persistence baseline:} The naive forecast $\hat{Y}(t) = Y(t-1)$
requires no training. For the delta formulation this is equivalent to
predicting $\hat{\Delta Y} = 0$ for every month. Persistence is a strong
benchmark because of high autocorrelation in conflict data: in always-active
high-conflict regions it is near-unbeatable on MAE alone. The CAST+
evaluation reports performance against both the RF baseline and the
persistence baseline.

\subsection{CAST+ Model: Two-Stage CatBoost}

The CAST+ model replaces single-stage regression with a two-stage
pipeline specifically designed for the zero-inflated delta distribution.

\textbf{Stage 1 (Classifier):} A \texttt{CatBoostClassifier} predicts
whether any change will occur in month $t$, that is whether
$\Delta Y \neq 0$. It is trained on the full training set with binary
target $\mathbf{1}[\Delta Y \neq 0]$. The change threshold is set at 0.50.

\textbf{Stage 2 (Regressor):} A \texttt{CatBoostRegressor} predicts the
signed magnitude of change, trained exclusively on nonzero-delta months.
This forces the regressor to specialise on the hard problem of estimating
how much conflict changes, rather than wasting capacity learning to
predict zero.

\textbf{Combined prediction:}
\[
    \hat{\Delta Y} = \text{Stage1\_flag} \times \text{Stage2\_magnitude}
\]

When Stage 1 predicts no change, the combined prediction is zero. When
Stage 1 predicts change, the combined prediction is the Stage 2 magnitude.

\textbf{Hyperparameters:} Both stages use 300 iterations, learning rate
0.05, depth 4, random seed 42. CatBoost's native categorical feature
handling is used for religion columns and income level code, with raw NaN
passthrough (filling NaN with zero was tested and found inferior).

\textbf{Per-target lag strategy:} The lag structure differs by target
based on empirical results from diagnostic experiments. Battles uses
t-1 lags only, as it is the most volatile target and extended lags
introduce noise. Remote violence and Violence against civilians use
t-1, t-6, and t-12 lags, as they benefit from seasonal signals without
the noise penalty observed for Battles.

\textbf{Sample weighting:} Exponential recency weights are applied during
training, down-weighting older observations to account for structural
changes in conflict dynamics over time.

\subsection{Alternative Model: LGBM-Tweedie with Country Hierarchy}

As a complementary model stream, a LightGBM regressor with a Tweedie
objective ($p = 1.5$) was evaluated alongside a systematic feature
ablation study. The Tweedie distribution models a compound
Poisson-Gamma process suited to non-negative, zero-inflated, and
overdispersed conflict counts---a better fit than Poisson (which
assumes mean equals variance) or squared-error objectives (which
penalise large counts symmetrically). Hyperparameters:
$n_\text{estimators} = 200$, learning rate $0.05$,
\texttt{num\_leaves} $= 31$, $\alpha = 0.1$, $\lambda = 1.0$.
Exponential recency weights ($w = \exp(-0.05 \times \text{months
since latest})$) are applied during training.

The ablation study evaluated five model families (Random Forest,
LightGBM-Poisson, LightGBM-Tweedie, XGBoost, Gradient Boosted
Regression) across six cumulative feature tiers (Baseline, +Risk,
+Macro, +Holidays, +Country, +Engineered) on the ten most
historically active conflict regions and three regression targets,
yielding 900 evaluations (V2). The defining result is a
discontinuous $-4.54$ MAE ($-7.9\%$) improvement from the +Country
(leave-one-out hierarchy) tier for LGBM-Tweedie---the single largest
gain across all tiers and model families. Full ablation results are
presented in Section~6.4.


% ============================================================
\section{Results}

\subsection{Reference Points}

Before presenting CAST+ results, two reference points are established.

The \textbf{CAST baseline} is a Random Forest trained on ACLED-only
features (29 predictors) without any enrichment. This is the starting
point for all comparisons. On the global evaluation (2,205 regions),
the CAST baseline achieves an overall MAE of 2.344. On the top-10 active
regions, it achieves an overall MAE of 67.968.

The \textbf{persistence baseline} predicts no change each month
($\hat{\Delta Y} = 0$). On the global evaluation, persistence achieves
an overall MAE of 0.931. This is lower than the CAST baseline because
86 to 92\% of months genuinely have zero change, so predicting zero is
often correct on MAE. However, persistence by definition achieves zero
directional accuracy for escalating and de-escalating months.

\subsection{CAST+ Global Results (Primary)}

Table~\ref{tab:final_global} presents the primary evaluation of CAST+
against the CAST baseline across all 2,205 valid Admin-1 regions.

\begin{table}[h]
\centering
\caption{Final comparison: CAST baseline vs CAST+ on global evaluation
(2,205 regions, 6-month holdout).}
\label{tab:final_global}
\begin{tabular}{lrrrrrr}
\toprule
Target & CAST MAE & CAST+ MAE & Improvement & CAST MAPE & CAST+ MAPE & Dir. Acc. \\
\midrule
Battles            & 1.782 & 1.384 & +22.3\% & 96.3\%  & 96.1\%  & 78.9\% \\
Remote Violence    & 4.319 & 2.546 & +41.0\% & 125.3\% & 93.1\%  & 81.6\% \\
VAC                & 0.932 & 0.836 & +10.3\% & 79.1\%  & 79.3\%  & 79.0\% \\
\midrule
\textbf{Overall}   & \textbf{2.344} & \textbf{1.589} & \textbf{+24.6\%} & & & \\
\bottomrule
\end{tabular}
\end{table}

CAST+ achieves a mean MAE improvement of 24.6\% over the CAST baseline,
surpassing the project target of 20\%. The largest gain is for Remote
violence (+41.0\%), where the delta formulation and two-stage
architecture most effectively separate the change signal from the
persistent zero-heavy baseline. Directional accuracy of 79 to 82\%
across all three targets confirms the model is learning meaningful
directional patterns rather than simply copying recent conflict levels.

The overall progression from baseline to CAST+ is summarised in
Table~\ref{tab:progression}.

\begin{table}[h]
\centering
\caption{Overall performance progression.}
\label{tab:progression}
\begin{tabular}{llrrrl}
\toprule
Stage & Model & Battles & Remote Viol. & VAC & Overall MAE \\
\midrule
CAST Baseline & RF (ACLED only)    & 1.782 & 4.319 & 0.932 & 2.344 \\
Persistence   & $Y(t) = Y(t-1)$   & 0.902 & 1.119 & 0.771 & 0.931 \\
\textbf{CAST+} & \textbf{Two-Stage CatBoost} & \textbf{1.384} & \textbf{2.546} & \textbf{0.836} & \textbf{1.589} \\
\bottomrule
\end{tabular}
\end{table}

\subsection{CAST+ vs CAST Baseline on Top-10 Active Regions}

Table~\ref{tab:top10_comparison} presents the fair comparison on the
top-10 most active regions, using the same RF ACLED-only baseline.

\begin{table}[h]
\centering
\caption{Fair comparison on top-10 most active regions.}
\label{tab:top10_comparison}
\begin{tabular}{lrrr}
\toprule
Target & CAST Baseline MAE & CAST+ MAE & Improvement \\
\midrule
Battles         & 63.938 & 29.615  & +53.7\% \\
Remote Violence & 127.930 & 120.414 & +5.9\%  \\
VAC             & 12.036 & 6.842   & +43.2\% \\
\midrule
\textbf{Overall} & \textbf{67.968} & \textbf{52.290} & \textbf{+23.1\%} \\
\bottomrule
\end{tabular}
\end{table}

CAST+ achieves a 23.1\% improvement over the CAST baseline on the
top-10 most active regions, and an overall MAE of 52.290 compared to
the CAST baseline of 67.968. CAST+ dominates on Battles (+53.7\%) and
VAC (+43.2\%). Remote violence shows a smaller improvement (+5.9\%),
reflecting the extreme scale of activity in the Ukraine conflict regions
during the holdout period.

\subsection{LGBM-Tweedie Ablation Results}

Tables~\ref{tab:v1_ablation} and \ref{tab:v2_ablation} report overall
MAE (mean across all three targets and the ten historical top-10
regions) for the V1 and V2 ablations. Feature sets accumulate
left to right.

\begin{table}[h]
\centering
\caption{V1 ablation: overall MAE by model and feature set
(historical top-10, 750 evaluations).}
\label{tab:v1_ablation}
\begin{tabular}{lrrrrr}
\toprule
\textbf{Model} & \textbf{Baseline} & \textbf{+Risk} & \textbf{+Macro} & \textbf{+Holidays} & \textbf{+Engineered} \\
\midrule
GBR                & 60.09 & 58.85 & 58.13 & 58.62 & 61.42 \\
LGBM-Poisson       & 53.63 & 53.73 & 53.65 & 53.69 & \textbf{53.20} \\
LGBM-Tweedie       & 57.48 & 57.67 & 57.27 & 57.36 & 57.73 \\
RF (CAST Baseline) & 57.25 & 57.95 & 56.64 & 57.12 & 57.37 \\
XGBoost            & 62.45 & 62.14 & 62.20 & 62.38 & 60.41 \\
\bottomrule
\end{tabular}
\end{table}

\begin{table}[h]
\centering
\caption{V2 ablation: overall MAE by model and feature set
(historical top-10, 900 evaluations). +Country tier added.}
\label{tab:v2_ablation}
\begin{tabular}{lrrrrrr}
\toprule
\textbf{Model} & \textbf{Baseline} & \textbf{+Risk} & \textbf{+Macro} & \textbf{+Holidays} & \textbf{+Country} & \textbf{+Engineered} \\
\midrule
GBR                   & 64.25 & 65.17 & 67.00 & 65.90 & 69.47          & 63.73 \\
LGBM-Poisson          & 55.64 & 55.84 & 55.78 & 55.85 & 56.83          & 56.71 \\
\textbf{LGBM-Tweedie} & 56.84 & 56.94 & 57.62 & 57.59 & \textbf{53.05} & 53.60 \\
RF (CAST Baseline)    & 58.98 & 59.10 & 60.48 & 59.43 & 60.11          & 60.08 \\
XGBoost               & 63.84 & 63.22 & 63.17 & 63.40 & 64.10          & 60.56 \\
\bottomrule
\end{tabular}
\end{table}

The +Country tier accounts for all of the V2 improvement for
LGBM-Tweedie ($-4.54$ MAE, $-7.9\%$); no other tier achieves more
than a 1.2\% change in either direction.

To enable a direct region-matched comparison with CAST+, the full
V2 ablation and the RF Baseline were re-run on the paper's
holdout-period top-10 regions (Table~\ref{tab:paper_top10_lgbm}).
LGBM-Tweedie +Country achieves MAE 73.06 on these regions, which is
5.2\% above the RF Baseline (69.47) and 40\% above CAST+ (52.29).
The reversal relative to the historical top-10 result ($-7.3\%$
vs RF) reflects the composition of the holdout set: seven of ten
regions are Ukraine frontlines, where the extreme scale and
volatility of the Russia-Ukraine war favours the delta+two-stage
architecture over an absolute-count regressor.

\begin{table}[h]
\centering
\caption{Performance on the paper's holdout-period top-10 regions.
DA-nonzero = directional accuracy on months where conflict changes.}
\label{tab:paper_top10_lgbm}
\begin{tabular}{lrrrrl}
\toprule
\textbf{Model} & \textbf{Battles} & \textbf{Remote} & \textbf{VaC} & \textbf{Overall MAE} & \textbf{DA-nonzero} \\
\midrule
RF Baseline (ACLED-only)   & 64.34 & 132.56 & 11.51 & \textbf{69.47} & --- \\
LGBM-Tweedie +Country      & 66.68 & 140.21 & 12.29 & \textbf{73.06} & 0.65 \\
CAST+ (Two-Stage CatBoost) & 29.62 & 120.41 &  6.84 & \textbf{52.29} & 0.79--0.82 \\
\bottomrule
\end{tabular}
\end{table}

The global LGBM-Tweedie evaluation covers 2,262 regions (requiring
at least 12 months of data and non-zero training target activity),
compared to 2,205 for CAST+ (which additionally requires variation
in monthly delta targets). The 57-region difference reflects the
different validity criteria appropriate to each model's objective
and does not indicate a methodological inconsistency. On the metric
of directional accuracy restricted to changing months (DA-nonzero),
LGBM-Tweedie achieves 0.84 globally---directly comparable to
CAST+'s 79--82\%, and substantially higher than the raw all-months
DA of 0.21, which is depressed by the Tweedie objective's structural
inability to predict zero in the 83\% of region-months where
conflict does not change.

\subsection{Conflict Intensity Segmentation}

To understand where CAST+ adds most value, regions were split into three
tiers per target based on mean monthly event counts: quiet (less than 1),
moderate (1 to 10), and active (more than 10).

Key findings from this analysis are as follows. Active Remote violence
regions show a 92\% MAE improvement over persistence, the strongest
result in the entire evaluation. This is where the two-stage architecture
adds most value, correctly identifying escalation events in regions with
established but volatile conflict patterns. Quiet Battles regions show
CAST+ winning on MAE, indicating usefulness for early conflict onset
detection. For moderate and active Battles and VAC tiers, persistence
is competitive on MAE, but CAST+ wins on directional accuracy,
demonstrating that even where the absolute error is comparable, CAST+
better identifies the direction of change.


% ============================================================
\section{Discussion}

\subsection{What Worked}

\textbf{Architecture was the biggest driver of improvement.} The
two-stage design directly addresses the zero-inflation problem. Separating
the question of whether change occurs from how much it changes allows
each stage to specialise. Stage 1 achieves meaningful precision on
identifying escalating months. Stage 2, trained only on nonzero-delta
months, learns the magnitude structure without the distortion of
zero-heavy training data.

\textbf{Delta formulation is the right approach.} Predicting change and
recovering counts is conceptually cleaner and empirically better than
predicting raw counts directly. It decouples the structural level
(captured by $Y(t-1)$) from the month-to-month signal (captured by the
model), making the model's task tractable.

\textbf{Per-target lag strategy matters.} Remote violence and VAC
benefit from seasonal lags (t-6, t-12) while Battles does not. A
one-size-fits-all lag structure leaves performance on the table.

\textbf{World Bank features contribute structurally.} Income inequality
and youth unemployment encode background vulnerability conditions that
complement the autoregressive conflict signal, particularly for the
Stage 1 classifier.

\textbf{Holiday and religion interaction features are principled.}
Although their marginal impact is modest, contextualising holiday counts
by religious composition is conceptually correct and contributes to
Stage 1 at low but nonzero importance.

\subsection{What Did Not Work}

\textbf{Risk features add noise in practice.} Despite theoretical
relevance, removing all risk indicator features improves MAE across
all three targets in the ablation study. These features are candidates
for removal in a production pipeline.

\textbf{Religion features alone have no contribution.} Without
interaction with holiday timing, static religious composition data
provides no predictive signal.

\textbf{Country hierarchy features provided minimal gain} (less than
1\%) in the two-stage CatBoost framework, despite being beneficial
in the LGBM-Tweedie context. This reflects the different way in which
the two architectures handle the zero-inflation problem.

\textbf{MAPE is not appropriate for this data.} The 86 to 92\%
zero-delta months cause division-by-zero instability in MAPE
calculations. Directional accuracy is a more honest metric for
evaluating 6-month-ahead conflict forecasting.

\subsection{Why Persistence Remains Competitive}

Persistence achieves lower overall MAE than both the CAST baseline
and CAST+ on the global evaluation (0.931 vs 1.589 and 2.344
respectively). This is not a failure of the model but a property of
the data. When 86 to 92\% of delta values are zero, predicting no
change is trivially accurate on MAE. Persistence has zero ability to
detect escalation or de-escalation by definition.

CAST+ meaningfully beats persistence on the metrics that matter for
early warning: directional accuracy (79 to 82\% vs 50\% for a random
classifier) and MAE in active Remote violence regions (92\%
improvement). The aggregate MAE comparison is not the right lens for
evaluating an early warning system.

\subsection{Why the Two Evaluation Setups Differ}

The global evaluation (2,205 regions) and the top-10 evaluation are
complementary. The global evaluation measures overall system
performance across the full distribution of conflict intensity, which
is the appropriate metric for a global forecasting tool. The top-10
evaluation measures performance in the world's most active conflict
zones during the holdout period, which is dominated by the Ukraine
war (7 of the top 10 regions). In sustained peak-intensity war zones,
all models including the CAST baseline show compressed performance
differences because the signal environment is qualitatively different
from the full global distribution.

\subsection{Potential Future Direction: Hybrid Per-Target Architecture}

The results across both model streams suggest a potential further
improvement: a hybrid architecture where the best-performing model
is selected per target rather than using a single model for all three.
Specifically, two-stage CatBoost delivers the strongest performance
on Battles (+53.7\% on top-10, +22.3\% globally). On the
holdout-period top-10 evaluation, CAST+ also outperforms
LGBM-Tweedie for Remote violence (120.41 vs 138.37) and Violence
against civilians (6.84 vs 12.29), so no single alternative model
provides an advantage on any individual target in this evaluation
set. A per-target hybrid architecture remains a direction worth
exploring as the evaluation expands beyond the
Ukraine-dominated holdout period.


% ============================================================
\section{Limitations and Future Work}

\textbf{Single holdout window.} All results are based on a single
six-month chronological holdout. Walk-forward rolling-origin validation
would provide more stable MAE estimates and reduce sensitivity to the
specific events of the holdout period.

\textbf{t-1 dependency.} The model heavily relies on t-1 conflict lags
as its strongest predictors. For a 6-month-ahead forecasting tool this
creates a practical limitation: the most recent month's data is
available when the forecast is made, but the model's performance on
longer horizons within the six-month window has not been separately
evaluated.

\textbf{Structural breaks.} The model will likely underperform on
structural breaks, defined as sudden onset or cessation of conflict
with no historical precedent. Conflict lags and momentum features
cannot anticipate genuinely novel escalation events. This is a
fundamental property of any autoregressive forecasting system.

\textbf{Risk features and LLM quality.} The news-scraped risk features
add noise rather than signal in the current implementation. This may
reflect data quality limitations in the LLM-parsed outputs, the
country-to-Admin-1 broadcast losing granularity, or temporal mismatch
between news mentions and actual events.

\textbf{Religion and holiday features.} While the interaction features
are conceptually sound, their contribution is modest. More granular
subnational religious composition data, or event-driven religious
tension signals, could improve their predictive value.

\textbf{Broader evaluation.} The top-10 evaluation subset is dominated
by Ukraine (7 of 10 regions during the holdout period). Expanding to
top-50 or top-100 would include more diverse conflict types where
trajectory features and structural signals may provide larger gains.

\textbf{Hierarchical and graph-based models.} The country hierarchy
features tested here are a tree-compatible approximation of a
hierarchical prior. A proper Bayesian mixed-effects model with country
as a random intercept, or a graph neural network using the Admin-1
spatial adjacency graph, would additionally share statistical strength
across low-activity regions, propagate cross-border spillover through
the full graph, and learn country-specific trend slopes.


% ============================================================
\section{Conclusion}

This report presents CAST+, an augmented conflict forecasting pipeline
that improves on the ACLED CAST baseline through enriched data
integration and a novel two-stage model architecture. The two-stage
CatBoost model, operating on monthly conflict deltas with per-target
lag strategies and native categorical feature handling, achieves a
24.6\% improvement in mean absolute error over the CAST baseline on a
global evaluation across 2,205 Admin-1 regions, surpassing the project
target of 20\%. On the top-10 most active regions, CAST+ achieves a
23.1\% improvement with an overall MAE of 52.290, outperforming all
alternative model configurations tested. Directional accuracy of 79 to
82\% confirms the model is detecting meaningful escalation and
de-escalation patterns.

The core methodological contributions are the delta formulation of the
prediction target, the two-stage zero-inflation-aware architecture, the
per-target lag strategy, and the holiday-religion interaction features
that contextualise calendar effects within the religious structure of
each country. The key empirical finding is that architecture matters
more than feature enrichment: switching from single-stage regression
to the two-stage design is the single largest driver of improvement,
while most individual feature groups show minimal marginal gains.

The LGBM-Tweedie ablation study establishes that the country-level
leave-one-out hierarchy feature is the single most impactful
enrichment for absolute-count conflict regressors ($-7.9\%$ MAE in
isolation), and confirms that both model streams achieve comparable
directional accuracy ($\approx 80$--$84\%$) when evaluated on the
same basis of changing months only. The two streams are structurally
complementary: CAST+'s two-stage delta architecture excels at change
detection across the full global distribution, while the LGBM-Tweedie
ablation provides a cross-model feature map identifying which
enrichments generalise beyond a single architecture.

All code, data pipelines, and evaluation notebooks are available at
\url{https://github.com/rapsoj/predicting-global-conflict} (Rosa-Branch).


% ============================================================
\appendix

\section{Repository Structure}

The full project code is available at
\url{https://github.com/rapsoj/predicting-global-conflict} on the
Rosa-Branch. Key files for replicating the results in this report are:

\begin{itemize}
    \item \texttt{forecast\_model/notebooks/final\_rosa.ipynb}: primary
          CAST+ two-stage CatBoost evaluation.
    \item \texttt{forecast\_model/notebooks/change\_target\_track.ipynb}:
          delta formulation experiments, lag strategy ablation,
          intensity segmentation.
    \item \texttt{forecast\_model/notebooks/model\_comparison\_4.ipynb}:
          multi-model ablation benchmark (RF, LGBM-Poisson,
          LGBM-Tweedie, XGBoost, GBR).
    \item \texttt{forecast\_model/utils/features/}: World Bank, holiday,
          and religion feature engineering modules.
    \item \texttt{forecast\_model/notebooks/WBD\_features.py}: engineered
          World Bank derived features.
    \item \texttt{forecast\_model/notebooks/hol\_rel.py}: holiday-religion
          interaction features.
    \item \texttt{src/scraping/}: news scraping pipeline.
    \item \texttt{forecast\_model/utils/risk\_merge.py}: risk indicator
          merge methodology.
\end{itemize}

\section{Feature Reference Table}

All features in the enriched dataset. Y = used by model; N = not
used. Layers 5b and 7 are Rosa's additions not in
\texttt{model\_data\_v2\_enriched.csv}; Layers 8--9 are Giray's
additions not used by CAST+.

\begin{longtable}{p{6.2cm} c p{3cm} l c c}
\toprule
\textbf{Feature} & \textbf{Layer} & \textbf{Source} & \textbf{Lag} & \textbf{LGBM} & \textbf{CAST+} \\
\midrule
\endfirsthead
\multicolumn{6}{c}{\textit{(continued)}} \\
\toprule
\textbf{Feature} & \textbf{Layer} & \textbf{Source} & \textbf{Lag} & \textbf{LGBM} & \textbf{CAST+} \\
\midrule
\endhead
\midrule \multicolumn{6}{r}{\textit{continued on next page}} \\
\endfoot
\bottomrule
\endlastfoot
\multicolumn{6}{l}{\textit{Layer 1 — ACLED Autoregressive}} \\
Battles $(t-1)$ & 1 & ACLED & $t-1$ & Y & Y \\
Explosions/Remote violence $(t-1)$ & 1 & ACLED & $t-1$ & Y & Y \\
Protests $(t-1)$ & 1 & ACLED & $t-1$ & Y & Y \\
Riots $(t-1)$ & 1 & ACLED & $t-1$ & Y & Y \\
Strategic developments $(t-1)$ & 1 & ACLED & $t-1$ & Y & Y \\
Violence against civilians $(t-1)$ & 1 & ACLED & $t-1$ & Y & Y \\
Excessive force against protesters $(t-1)$ & 1 & ACLED & $t-1$ & Y & Y \\
Agreement $(t-1)$ & 1 & ACLED & $t-1$ & Y & Y \\
Battles\_neighbours $(t-1)$ & 1 & ACLED spatial & $t-1$ & Y & Y \\
Remote\_neighbours $(t-1)$ & 1 & ACLED spatial & $t-1$ & Y & Y \\
Protests\_neighbours $(t-1)$ & 1 & ACLED spatial & $t-1$ & Y & Y \\
Riots\_neighbours $(t-1)$ & 1 & ACLED spatial & $t-1$ & Y & Y \\
Strat.\ dev.\_neighbours $(t-1)$ & 1 & ACLED spatial & $t-1$ & Y & Y \\
VaC\_neighbours $(t-1)$ & 1 & ACLED spatial & $t-1$ & Y & Y \\
\multicolumn{6}{l}{\textit{Layer 2 — Temporal}} \\
linear\_month\_trend & 2 & Derived & none & Y & Y \\
year & 2 & Derived & none & Y & Y \\
month\_sin, month\_cos & 2 & Derived (cyclic) & none & Y & Y \\
quarter\_sin, quarter\_cos & 2 & Derived (cyclic) & none & Y & Y \\
\multicolumn{6}{l}{\textit{Layer 3 — News-Scraped Risk Indicators}} \\
risk\_contested\_election $(t-1)$ & 3 & GNews/LLM & $t-1$ & Y & Y \\
risk\_crop\_failure $(t-1)$ & 3 & GNews/LLM & $t-1$ & Y & Y \\
risk\_economic\_concern $(t-1)$ & 3 & GNews/LLM & $t-1$ & Y & Y \\
risk\_ethnic\_tension $(t-1)$ & 3 & GNews/LLM & $t-1$ & Y & Y \\
risk\_military\_coup $(t-1)$ & 3 & GNews/LLM & $t-1$ & Y & Y \\
risk\_natural\_disaster $(t-1)$ & 3 & GNews/LLM & $t-1$ & Y & Y \\
risk\_political\_assassination $(t-1)$ & 3 & GNews/LLM & $t-1$ & Y & Y \\
\multicolumn{6}{l}{\textit{Layer 4 — World Bank Macroeconomic}} \\
inflation\_py & 4 & World Bank (CPI) & prior-year & Y & Y \\
youth\_unemployment\_py & 4 & World Bank & prior-year & Y & Y \\
income\_inequality\_py & 4 & World Bank (Gini) & prior-year & Y & Y \\
income\_level\_code & 4 & World Bank metadata & structural & Y & Y \\
\multicolumn{6}{l}{\textit{Layer 5a — Religious Composition (WRP 2010)}} \\
majority\_religion, majority\_pct & 5a & WRP & static & Y & Y \\
minority1\_religion, minority1\_pct & 5a & WRP & static & Y & Y \\
minority2\_religion, minority2\_pct & 5a & WRP & static & Y & Y \\
nonreligpct & 5a & WRP & static & Y & Y \\
\multicolumn{6}{l}{\textit{Layer 5b — Engineered World Bank Features (Rosa)}} \\
structural\_vulnerability & 5b & Derived (WB) & prior-year & N & Y \\
inflation\_shock & 5b & Derived (WB) & prior-year & N & Y \\
shock\_vulnerability & 5b & Derived (WB) & prior-year & N & Y \\
capacity\_adjusted\_risk & 5b & Derived (WB) & prior-year & N & Y \\
inflation\_change & 5b & Derived (WB) & prior-year & N & Y \\
high\_inequality\_flag & 5b & Derived (WB) & prior-year & N & Y \\
inflation\_shock\_poor & 5b & Derived (WB) & prior-year & N & Y \\
conflict\_momentum & 5b & Derived (ACLED) & $t-1$ & N & Y \\
violence\_escalation & 5b & Derived (ACLED) & $t-1$ & N & Y \\
macro\_conflict\_pressure & 5b & Derived (ACLED+WB) & $t-1$ & N & Y \\
\multicolumn{6}{l}{\textit{Layer 6 — Religion-Specific Holiday Features}} \\
holiday\_count\_month, is\_holiday\_month & 6 & Holiday calendar & current & Y & Y \\
christian/islam/shia/hindu/buddhist/jewish/ cultural/nonreligious\_holiday\_count & 6 & Holiday calendar & current & Y & Y \\
\multicolumn{6}{l}{\textit{Layer 7 — Holiday $\times$ Religion Interactions (Rosa)}} \\
weighted\_holidays\_majority/minority1/minority2 & 7 & Derived (holiday$\times$WRP) & $t-1$ & N & Y \\
minority\_tension\_minority1/minority2/total & 7 & Derived (holiday$\times$WRP) & $t-1$ & N & Y \\
total\_religious\_mobilization & 7 & Derived (holiday$\times$WRP) & $t-1$ & N & Y \\
\multicolumn{6}{l}{\textit{Layer 8 — Country-Level Hierarchy Features (Giray)}} \\
country\_battles\_excl $(t-1)$ & 8 & Derived (ACLED LOO) & $t-1$ & Y & N \\
country\_remote\_excl $(t-1)$ & 8 & Derived (ACLED LOO) & $t-1$ & Y & N \\
country\_vac\_excl $(t-1)$ & 8 & Derived (ACLED LOO) & $t-1$ & Y & N \\
country\_total\_excl $(t-1)$ & 8 & Derived (ACLED LOO) & $t-1$ & Y & N \\
\multicolumn{6}{l}{\textit{Layer 9 — Engineered Cross-Variable Features (Giray)}} \\
Battles/Remote/VaC $(t-2)$ & 9 & ACLED & $t-2$ & Y & N \\
organized\_violence $(t-1)$ & 9 & Derived (ACLED) & $t-1$ & Y & N \\
is\_active $(t-1)$ & 9 & Derived (ACLED) & $t-1$ & Y & N \\
battles\_x\_remote $(t-1)$ & 9 & Derived (ACLED) & $t-1$ & Y & N \\
Battles/Remote/VaC\_3mo\_avg $(t-1)$ & 9 & Derived (ACLED) & $t-1$ & Y & N \\
\end{longtable}

\end{document}